% 09_conclusion.tex -- the rehauled paper's section, as bullets to expand into prose (2026-10-15).
% Sources and pointers behind every number: writeup/PAPER_REHAUL_OUTLINE_2026-10-15.md, section 9.
% House rules: every number in the final prose comes from a generated macro (writeup/generated/*.tex), never
% typed; park (comment out with a dated note) rather than delete.
% Macros to use here: reuse the introduction's macros; no new numbers
% 2026-10-15 (notation pass, writeup/NOTATION_GUIDE_2026-10-15.md): symbols and wording; the typed 866 -> \cmNdays.
% Old -> new: RV_t -> \RV_t; B_t -> s^2_{t,c}; IV^{30}_t -> \mathrm{IV}_t; HAR + calendar -> HAR benchmark;
% exogenous inputs -> exogenous predictors; per bar / pooled -> closing-interval / pooled intraday model;
% Sharpe -> \SR; trade -> straddle strategy; always-short -> short-straddle benchmark; crossing the spread -> at
% quoted bid/ask prices. The guarded \input is a no-op after the introduction's.
% 2026-10-15 (consolidation): window row s -> t' (s is the periodicity factor); the training loss written with the
% realization floor and log-ratio clamp of Methods (FLAG (3) there).
% 2026-10-15 (consolidation): parked -- the number files are loaded once in the preamble of paper_rehaul.tex.
% \ifdefined\cmNdays\else% AUTO-GENERATED by writeup/make_table_close_main.py -- do not edit.
% Numbers quoted in the prose of writeup/sections/results_close_option.tex; each is read
% from results/close_master_table/master_table.csv and the sources in the script.
% 2026-10-06: recalibration removed / intervals removed, user decision -- the plain forecast
% f^2 B, point estimates; ...T = HAC t of the daily mid-return difference, ...DM = Diebold-
% Mariano statistic; a ...Above / ...Below count counts point estimates, its ...Sig twin those
% significant at 5 % (DM for QLIKE, HAC t for Sharpe).
\newcommand{\cmNdays}{866}
\newcommand{\cmNrows}{22}
\newcommand{\cmShortMean}{0.014}
\newcommand{\cmShortShMid}{0.20}
\newcommand{\cmShortShX}{\ensuremath{-}0.27}
\newcommand{\cmShortT}{\ensuremath{-}1.71}
\newcommand{\cmPaperShMin}{0.43}
\newcommand{\cmPaperShMax}{1.50}
\newcommand{\cmPaperXMin}{\ensuremath{-}0.05}
\newcommand{\cmPaperXMax}{1.01}
\newcommand{\cmPaperVsShortMin}{+0.22}
\newcommand{\cmPaperVsShortMax}{+1.30}
\newcommand{\cmPaperVsShortAboveN}{8}
\newcommand{\cmPaperVsShortAboveSigN}{three}
\newcommand{\cmPaperVsShortAboveSigWho}{LightGBM, XGBoost, lasso (causally tuned)}
\newcommand{\cmPaperSameMin}{82}
\newcommand{\cmSellPaperMin}{72}
\newcommand{\cmSellPaperMax}{83}
\newcommand{\cmPaperSameMax}{96}
\newcommand{\cmRefQ}{0.1092}
\newcommand{\cmRefShMid}{1.28}
\newcommand{\cmRefShX}{0.80}
\newcommand{\cmRefBuy}{27}
\newcommand{\cmRefVsShort}{+1.08}
\newcommand{\cmRefVsShortT}{+1.71}
\newcommand{\cmFomcSh}{1.36}
\newcommand{\cmFomcD}{+0.08}
\newcommand{\cmFomcDT}{+0.44}
\newcommand{\cmLassoFShMid}{1.06}
\newcommand{\cmLassoFVsRef}{\ensuremath{-}0.22}
\newcommand{\cmLassoFVsRefT}{\ensuremath{-}0.75}
\newcommand{\cmLassoFSame}{96}
\newcommand{\cmPaperFarLabel}{baseline (HAR + calendar OLS)}
\newcommand{\cmPaperFarVsRef}{\ensuremath{-}0.85}
\newcommand{\cmPaperFarVsRefT}{\ensuremath{-}1.49}
\newcommand{\cmPaperBestLabel}{lasso (causally tuned)}
\newcommand{\cmPaperBestSh}{1.50}
\newcommand{\cmNlFamN}{five}
\newcommand{\cmNlMasterN}{90}
\newcommand{\cmPoolShMin}{0.91}
\newcommand{\cmPoolShMax}{1.44}
\newcommand{\cmSubShMin}{1.37}
\newcommand{\cmSubShMax}{1.92}
\newcommand{\cmNlShMin}{1.23}
\newcommand{\cmNlShMax}{1.81}
\newcommand{\cmTabVsRefN}{21}
\newcommand{\cmTabVsRefAbove}{15}
\newcommand{\cmTabVsRefBelow}{6}
\newcommand{\cmTabVsRefAboveSig}{0}
\newcommand{\cmTabVsRefBelowSig}{0}
\newcommand{\cmMasterN}{112}
\newcommand{\cmMasterNPaper}{9}
\newcommand{\cmMasterNBase}{52}
\newcommand{\cmMasterMedQBase}{0.1033}
\newcommand{\cmMasterMedShBase}{0.73}
\newcommand{\cmMasterMedShXBase}{0.25}
\newcommand{\cmMasterNAll}{51}
\newcommand{\cmMasterMedQAll}{0.1014}
\newcommand{\cmMasterMedShAll}{1.27}
\newcommand{\cmMasterMedShXAll}{0.79}
\newcommand{\cmMasterSetsN}{two}
\newcommand{\cmMasterVsRefN}{111}
\newcommand{\cmMasterVsRefAbove}{32}
\newcommand{\cmMasterVsRefBelow}{79}
\newcommand{\cmMasterVsRefAboveSig}{0}
\newcommand{\cmMasterVsRefBelowSig}{0}
\newcommand{\cmMasterVsShortAbove}{112}
\newcommand{\cmMasterVsShortBelow}{0}
\newcommand{\cmMasterVsShortAboveSig}{20}
\newcommand{\cmMasterVsShortBelowSig}{0}
\newcommand{\cmMasterVsShortN}{112}
\newcommand{\cmMasterQVsRefBelow}{99}
\newcommand{\cmMasterQVsRefAbove}{12}
\newcommand{\cmMasterQVsRefBelowSig}{23}
\newcommand{\cmMasterQVsRefAboveSig}{1}
\newcommand{\cmMasterMedQ}{0.1028}
\newcommand{\cmMasterMedSh}{1.08}
\newcommand{\cmMasterMedShX}{0.59}
\newcommand{\cmMasterShMin}{0.29}
\newcommand{\cmMasterShMax}{1.92}
\newcommand{\cmMasterXMin}{\ensuremath{-}0.19}
\newcommand{\cmMasterXMax}{1.43}
\newcommand{\cmMasterVsShortMed}{+0.87}
\newcommand{\cmMasterRankCorr}{\ensuremath{-}0.20}
\newcommand{\cmMasterRankCorrX}{\ensuremath{-}0.20}
\newcommand{\cmSellMasterMin}{66}
\newcommand{\cmSellMasterMax}{83}
\newcommand{\cmSellMasterMed}{70}
\newcommand{\cmTopShLabel}{per-bar elastic net}
\newcommand{\cmTopShInputs}{\texttt{all\_features}}
\newcommand{\cmTopShMid}{1.92}
\newcommand{\cmXLossMin}{0.48}
\newcommand{\cmXLossMax}{0.50}
\newcommand{\cmLdSpec}{+0.16}
\newcommand{\cmLdSpecQ}{\ensuremath{-}4.3}
\newcommand{\cmLdSpecT}{+0.34}
\newcommand{\cmLdSpecDM}{\ensuremath{-}1.73}
\newcommand{\cmLdPerBarBase}{+0.46}
\newcommand{\cmLdPerBarBaseQ}{\ensuremath{-}12.0}
\newcommand{\cmLdPerBarBaseT}{+1.05}
\newcommand{\cmLdPerBarBaseDM}{\ensuremath{-}3.37}
\newcommand{\cmLdPerBarAll}{\ensuremath{-}0.03}
\newcommand{\cmLdPerBarAllQ}{\ensuremath{-}6.7}
\newcommand{\cmLdPerBarAllT}{\ensuremath{-}0.07}
\newcommand{\cmLdPerBarAllDM}{\ensuremath{-}1.61}
\newcommand{\cmLdIn}{+0.05}
\newcommand{\cmLdInQ}{+0.8}
\newcommand{\cmLdInT}{+0.09}
\newcommand{\cmLdInDM}{+0.18}
\newcommand{\cmLdEstMin}{\ensuremath{-}0.15}
\newcommand{\cmLdEstMax}{+0.50}
\newcommand{\cmLdEstMinLabel}{\texttt{subtree\_rf\_all\_features}}
\newcommand{\cmLdEstMaxLabel}{\texttt{subpt\_enet\_all\_features}}
\newcommand{\cmLdEstMinWho}{random forest}
\newcommand{\cmLdEstMaxWho}{elastic net}
\newcommand{\cmLdN}{21}
\newcommand{\cmLdNPath}{12}
\newcommand{\cmLdNCampaign}{9}
\newcommand{\cmLdResolved}{none}
\newcommand{\cmLdQResolved}{three}
\newcommand{\cmLdCDedup}{+0.00}
\newcommand{\cmLdCDedupQ}{+0.0}
\newcommand{\cmLdCMaskTten}{\ensuremath{-}0.21}
\newcommand{\cmLdCMaskTtenQ}{\ensuremath{-}1.0}
\newcommand{\cmLdCMaskTtenT}{\ensuremath{-}0.69}
\newcommand{\cmLdCMaskTtenDM}{\ensuremath{-}1.04}
\newcommand{\cmLdCMaskTone}{\ensuremath{-}0.23}
\newcommand{\cmLdCMaskToneQ}{+1.4}
\newcommand{\cmLdCMaskToneT}{\ensuremath{-}0.55}
\newcommand{\cmLdCMaskToneDM}{+1.36}
\newcommand{\cmLdCMaskRsten}{\ensuremath{-}0.08}
\newcommand{\cmLdCMaskRstenQ}{+2.8}
\newcommand{\cmLdCMaskRstenT}{\ensuremath{-}0.32}
\newcommand{\cmLdCMaskRstenDM}{+1.42}
\newcommand{\cmLdCMaskRsone}{\ensuremath{-}0.05}
\newcommand{\cmLdCMaskRsoneQ}{+0.2}
\newcommand{\cmLdCMaskRsoneT}{\ensuremath{-}0.17}
\newcommand{\cmLdCMaskRsoneDM}{+0.12}
\newcommand{\cmLdCCadLgbm}{\ensuremath{-}0.04}
\newcommand{\cmLdCCadLgbmQ}{\ensuremath{-}2.2}
\newcommand{\cmLdCCadLgbmT}{\ensuremath{-}0.08}
\newcommand{\cmLdCCadLgbmDM}{\ensuremath{-}1.12}
\newcommand{\cmLdCTuneRs}{+0.28}
\newcommand{\cmLdCTuneRsQ}{+3.0}
\newcommand{\cmLdCTuneRsT}{+0.58}
\newcommand{\cmLdCTuneRsDM}{+1.12}
\newcommand{\cmLdCTuneOp}{\ensuremath{-}0.10}
\newcommand{\cmLdCTuneOpQ}{+7.3}
\newcommand{\cmLdCTuneOpT}{\ensuremath{-}0.31}
\newcommand{\cmLdCTuneOpDM}{+2.83}
\newcommand{\cmLdCTuneOpPer}{\ensuremath{-}0.32}
\newcommand{\cmLdCTuneOpPerQ}{+2.5}
\newcommand{\cmLdCTuneOpPerT}{\ensuremath{-}0.65}
\newcommand{\cmLdCTuneOpPerDM}{+0.75}
\newcommand{\cmPNewBase}{22}
\newcommand{\cmPOldBase}{34}
\newcommand{\cmPNewAll}{628}
\newcommand{\cmPOldAll}{640}
\newcommand{\cmNDup}{12}
\newcommand{\cmDedupN}{37}
\newcommand{\cmDedupQMax}{0.05}
\newcommand{\cmDedupPosMax}{one}
\newcommand{\cmKeptBase}{17}
\newcommand{\cmKeptBaseMin}{17}
\newcommand{\cmKeptBaseMax}{17}
\newcommand{\cmKeptAll}{369}
\newcommand{\cmKeptAllMin}{360}
\newcommand{\cmKeptAllMax}{378}
\newcommand{\cmMaskN}{16}
\newcommand{\cmMaskQMed}{0.4}
\newcommand{\cmMaskQAbove}{10}
\newcommand{\cmMaskShBelow}{14}
\newcommand{\cmMaskQAboveSig}{none}
\newcommand{\cmMaskShBelowSig}{one}
\newcommand{\cmTrVsRidgeN}{48}
\newcommand{\cmTrSetsN}{two}
\newcommand{\cmMaskSetsN}{two}
\newcommand{\cmTrVsRidgeQAbove}{46}
\newcommand{\cmTrVsRidgeShBelow}{41}
\newcommand{\cmTrVsRidgeQAboveSig}{21}
\newcommand{\cmTrVsRidgeShBelowSig}{one}
\newcommand{\cmCadN}{six}
\newcommand{\cmCadQBelow}{five}
\newcommand{\cmCadQBelowSig}{none}
\newcommand{\cmCadShUp}{\ensuremath{-}0.43}
\newcommand{\cmCadShUpT}{\ensuremath{-}1.03}
\newcommand{\cmRsQAbove}{6}
\newcommand{\cmRsQAboveSig}{two}
\newcommand{\cmRsN}{6}
\newcommand{\cmOpN}{24}
\newcommand{\cmOpQAbove}{24}
\newcommand{\cmOpQAboveSig}{13}
\newcommand{\cmTpN}{18}
\newcommand{\cmTpQBelow}{12}
\newcommand{\cmTpQBelowSig}{one}
\newcommand{\cmBudN}{24}
\newcommand{\cmBudQAbove}{17}
\newcommand{\cmBudQAboveSig}{two}
\newcommand{\cmBudQBelow}{7}
\newcommand{\cmBudLastMin}{27}
\newcommand{\cmBudLastMax}{37}
\newcommand{\cmBudLastExch}{20}
\newcommand{\cmNlBestSh}{1.81}
\newcommand{\cmNlBestLabel}{tuned per-bar LightGBM}
\newcommand{\cmNlBestInputs}{\texttt{all\_features}}
\newcommand{\cmNlBestVsRef}{+0.53}
\newcommand{\cmNlBestVsRefT}{+1.04}
\newcommand{\cmNlVsRefAbove}{19}
\newcommand{\cmNlVsRefBelow}{71}
\newcommand{\cmNlVsRefAboveSig}{0}
\newcommand{\cmNlVsRefBelowSig}{0}
\newcommand{\cmClassChunks}{505}
\newcommand{\cmClassPts}{20}
\newcommand{\cmClassLgbmPick}{18}
\newcommand{\cmClassMaxRelMin}{4.5}
\newcommand{\cmClassMaxRelMax}{5.5}
\newcommand{\cmSvN}{31}
\newcommand{\cmSvSetN}{112}
\newcommand{\cmSvSetPairsN}{3472}
\newcommand{\cmSvSetAboveMid}{755}
\newcommand{\cmSvSetAboveMidSig}{29}
\newcommand{\cmSvSetStraddleFirstMid}{4}
\newcommand{\cmSvSetWingN}{672}
\newcommand{\cmSvSetWingMidBelow}{668}
\newcommand{\cmSvSetAboveX}{410}
\newcommand{\cmSvSetAboveXSig}{0}
\newcommand{\cmSvSetStraddleFirstX}{3}
\newcommand{\cmSvSetWingXBelow}{672}
\newcommand{\cmSvSetBestMidLabel}{call vertical, w50}
\newcommand{\cmSvSetBestMidMed}{+0.37}
\newcommand{\cmSvSetBestMidAboveN}{94}
\newcommand{\cmSvSetBestMidAboveSig}{0}
\newcommand{\cmSvCostStraddle}{2.9}
\newcommand{\cmSvCostIbfMin}{4.0}
\newcommand{\cmSvCostIbfMax}{6.1}
\newcommand{\cmMasterPosTotalN}{112}
\newcommand{\cmMasterTopTwentyShareMed}{99}
\newcommand{\cmMasterTopTwentyShareMin}{59}
\newcommand{\cmMasterTopTwentyShareMax}{338}
\newcommand{\cmMasterTopTwentyBuysN}{112}
\newcommand{\cmMasterBottomTwentySellsN}{112}
\newcommand{\cmMasterDiffMed}{+53.5}
\newcommand{\cmMasterDiffPosN}{112}
\newcommand{\cmMasterDiffExTenMed}{\ensuremath{-}30.0}
\newcommand{\cmMasterDiffExTenHalfN}{112}
\newcommand{\cmMasterDiffExTenNegN}{102}
\newcommand{\cmBContamMaxDSh}{0.38}
\newcommand{\cmBContamMaxDShX}{0.38}
\newcommand{\cmBContamMedDSh}{0.08}
\newcommand{\cmBContamPaperMaxDSh}{0.21}
\newcommand{\cmBContamPaperMaxDShX}{0.22}
\newcommand{\cmBContamPaperMedDSh}{0.08}
\newcommand{\cmBContamDays}{94}
\newcommand{\cmCensIvMaxDSh}{0.06}
\newcommand{\cmCensIvMaxDShX}{0.08}
\newcommand{\cmCensIvMedDSh}{0.04}
\newcommand{\cmCensIvPaperMaxDSh}{0.05}
\newcommand{\cmCensIvPaperMaxDShX}{0.07}
\newcommand{\cmCensIvPaperMedDSh}{0.04}
\newcommand{\cmCensIvDays}{5}
\newcommand{\cmBContamMaxFlips}{16}
\newcommand{\cmBContamMedFlips}{8}
\newcommand{\cmBContamPaperMaxFlips}{10}
\newcommand{\cmBContamMaxRel}{21}
\newcommand{\cmPtPreSessions}{3947}
\newcommand{\cmPtPreStart}{13~April~2004}
\newcommand{\cmPtValFirst}{6~January~2016}
\newcommand{\cmPtValLast}{2~January~2020}
\newcommand{\cmPtFolds}{four}
\newcommand{\cmPtFoldLen}{250}
\newcommand{\cmPtFitWin}{2{,}000}
\newcommand{\cmPtEmbargo}{25}
\newcommand{\cmPtRatio}{1/100}
\newcommand{\cmPtBlock}{250}
\newcommand{\cmPtRetuneN}{four}
\newcommand{\cmPtRetuneFirst}{23~June~2020}
\newcommand{\cmPtRetuneLast}{16~June~2023}
\newcommand{\cmPtFiredRidgeBase}{one}
\newcommand{\cmPtPenRidgeBase}{\ensuremath{10^{-2}}}
\newcommand{\cmPtFiredRidgeAll}{none}
\newcommand{\cmPtPenRidgeAll}{\ensuremath{10^{3}}}
\newcommand{\cmPtFiredLassoBase}{one}
\newcommand{\cmPtPenLassoBase}{\ensuremath{10^{-6}}}
\newcommand{\cmPtFiredLassoAll}{none}
\newcommand{\cmPtPenLassoAll}{\ensuremath{10^{-2}}}
\newcommand{\cmPtFiredEnetBase}{two}
\newcommand{\cmPtPenEnetBase}{\ensuremath{10^{-6}}}
\newcommand{\cmPtFiredEnetAll}{none}
\newcommand{\cmPtPenEnetAll}{\ensuremath{10^{-2}}}
\newcommand{\cmPtFiredTotal}{4}
\newcommand{\cmPtFiredOf}{24}
\newcommand{\cmPtCausalSwitches}{75}
\newcommand{\cmPtCausalChoices}{165}
\newcommand{\cmPtCausalMarginPass}{3}
\newcommand{\cmPtCausalEdgePicks}{132}
\newcommand{\cmPtCausalPicks}{198}
\newcommand{\cmPtSwitchFlipPct}{8}
\newcommand{\cmPtCausalSwitchesPaper}{24}
\newcommand{\cmPtCausalChoicesPaper}{45}
\newcommand{\cmPtCausalMarginPassPaper}{3}
\newcommand{\cmPtCausalEdgePicksPaper}{34}
\newcommand{\cmPtCausalPicksPaper}{54}
\newcommand{\cmPtSwitchFlipPctPaper}{9}
\newcommand{\cmPoolPtN}{6}
\newcommand{\cmPoolPtQPctMean}{+0.9}
\newcommand{\cmPoolPtQWorse}{5}
\newcommand{\cmPoolPtDMMax}{1.22}
\newcommand{\cmPoolPtTMax}{1.29}
\newcommand{\cmPoolPtSigN}{0}
\newcommand{\cmPoolPtShAllMin}{+0.17}
\newcommand{\cmPoolPtShAllMax}{+0.35}
\newcommand{\cmPoolPtShOtherMin}{\ensuremath{-}0.06}
\newcommand{\cmPoolPtShOtherMax}{\ensuremath{-}0.19}
\newcommand{\cmPoolPtBlocks}{218}
\newcommand{\cmPoolPtCausalSwMin}{137}
\newcommand{\cmPoolPtCausalSwMax}{170}
\newcommand{\cmPoolPtCausalMarginMin}{3}
\newcommand{\cmPoolPtCausalMarginMax}{22}
\newcommand{\cmPoolPtRetunes}{217}
\newcommand{\cmPoolPtRetuneSwMin}{11}
\newcommand{\cmPoolPtRetuneSwMax}{40}
\newcommand{\cmTreeProvRfName}{\texttt{rf\_opt\_00}}
\newcommand{\cmTreeProvRfTrials}{400}
\newcommand{\cmTreeProvRfBudget}{400}
\newcommand{\cmTreeProvRfDevMse}{0.108456}
\newcommand{\cmTreeProvLgbmDevMse}{0.107449}
\newcommand{\cmTreeProvXgbDevMse}{0.108076}
\newcommand{\cmTreeProvXgbTopDevMse}{0.107894}
\newcommand{\cmTreeProvRfTrees}{485}
\newcommand{\cmTreeProvRfDepth}{20}
\newcommand{\cmTreeProvRfFeatures}{0.90}
\newcommand{\cmTreeProvRfSamples}{0.41}
\newcommand{\cmTreeProvRfLeafDev}{12}
\newcommand{\cmTreeProvRfLeaf}{1}
\newcommand{\cmTreeProvLgbmLeafDev}{98}
\newcommand{\cmTreeProvLgbmLeaf}{8}
\newcommand{\cmTreeProvXgbTopName}{\texttt{xgb\_opt\_00}}
\newcommand{\cmTreeProvXgbName}{\texttt{xgb\_opt\_06}}
\newcommand{\cmTreeProvXgbDevRank}{7}
\newcommand{\cmTreeProvXgbArms}{8}
\newcommand{\cmTreeProvXgbPostRank}{8}
\newcommand{\cmTreeProvXgbRecentRank}{8}
\newcommand{\cmTreeProvXgbQRank}{8}
\newcommand{\cmTreeProvXgbPickDate}{14~August~2026}
\newcommand{\cmMasterNLean}{51}
\newcommand{\cmMasterMedQLean}{0.0991}
\newcommand{\cmMasterMedShLean}{1.37}
\newcommand{\cmMasterMedShXLean}{0.88}
\newcommand{\cmLdPerBarLean}{+0.65}
\newcommand{\cmLdPerBarLeanQ}{\ensuremath{-}12.3}
\newcommand{\cmLdPerBarLeanT}{+1.49}
\newcommand{\cmLdPerBarLeanDM}{\ensuremath{-}3.39}
\newcommand{\cmLdInBaseLean}{+0.51}
\newcommand{\cmLdInBaseLeanQ}{\ensuremath{-}3.3}
\newcommand{\cmLdInBaseLeanT}{+1.66}
\newcommand{\cmLdInBaseLeanDM}{\ensuremath{-}1.17}
\newcommand{\cmLdInLeanAll}{\ensuremath{-}0.46}
\newcommand{\cmLdInLeanAllQ}{+4.3}
\newcommand{\cmLdInLeanAllT}{\ensuremath{-}1.15}
\newcommand{\cmLdInLeanAllDM}{+1.50}
%
% 2026-10-06 parked (recalibration removed / intervals removed, user decision):
%   \cmShortShCiMid: interval of always short's Sharpe
%   \cmShortShCiX: interval of always short's Sharpe
%   \cmPaperVsShortAboveWho: the paper's forecasts whose interval vs always short lay above zero
%   \cmPaperVsShortAboveCi: that forecast's interval
%   \cmSellHead: the former headline's sell share (no headline model)
%   \cmHQ: the former headline (no headline model)
%   \cmHQPct: the former headline (no headline model)
%   \cmHQCi: the former headline (no headline model)
%   \cmHDM: the former headline (no headline model)
%   \cmHDMp: the former headline (no headline model)
%   \cmHShMid: the former headline (no headline model)
%   \cmHShX: the former headline (no headline model)
%   \cmHMean: the former headline (no headline model)
%   \cmHHit: the former headline (no headline model)
%   \cmHBuy: the former headline (no headline model)
%   \cmHVsRefMid: the former headline (no headline model)
%   \cmHVsRefX: the former headline (no headline model)
%   \cmHVsShortMid: the former headline (no headline model)
%   \cmHVsShortX: the former headline (no headline model)
%   \cmHVsRefT: the former headline (no headline model)
%   \cmHVsRefP: the former headline (no headline model)
%   \cmHVsShortT: the former headline (no headline model)
%   \cmHVsShortP: the former headline (no headline model)
%   \cmHRankSh: the former headline (no headline model)
%   \cmHRankQ: the former headline (no headline model)
%   \cmVsHN: the former headline (no headline model)
%   \cmVsHAbove: the former headline (no headline model)
%   \cmVsHBelow: the former headline (no headline model)
%   \cmVsHAboveSig: the former headline (no headline model)
%   \cmVsHBelowSig: the former headline (no headline model)
%   \cmVsHLassoSh: the former headline (no headline model)
%   \cmVsHLasso: paired against the former headline (no headline model)
%   \cmVsHLassoT: paired against the former headline (no headline model)
%   \cmVsHLassoSame: paired against the former headline (no headline model)
%   \cmVsHLassoQPct: paired against the former headline (no headline model)
%   \cmVsHLassoDM: paired against the former headline (no headline model)
%   \cmVsHEnet: paired against the former headline (no headline model)
%   \cmVsHEnetT: paired against the former headline (no headline model)
%   \cmVsHEnetSame: paired against the former headline (no headline model)
%   \cmVsHEnetQPct: paired against the former headline (no headline model)
%   \cmVsHEnetDM: paired against the former headline (no headline model)
%   \cmVsHRidgeAll: paired against the former headline (no headline model)
%   \cmVsHRidgeAllT: paired against the former headline (no headline model)
%   \cmVsHRidgeAllSame: paired against the former headline (no headline model)
%   \cmVsHRidgeAllQPct: paired against the former headline (no headline model)
%   \cmVsHRidgeAllDM: paired against the former headline (no headline model)
%   \cmVsHTop: paired against the former headline (no headline model)
%   \cmVsHTopT: paired against the former headline (no headline model)
%   \cmVsHTopSame: paired against the former headline (no headline model)
%   \cmVsHTopQPct: paired against the former headline (no headline model)
%   \cmVsHTopDM: paired against the former headline (no headline model)
%   \cmLdPerBarLive: a step to the former headline (no headline, no live-feasible)
%   \cmLdPerBarLiveQ: a step to the former headline (no headline, no live-feasible)
%   \cmLdPerBarLiveT: a step to the former headline (no headline, no live-feasible)
%   \cmLdPerBarLiveDM: a step to the former headline (no headline, no live-feasible)
%   \cmLdInAll: a step to the former headline (no headline, no live-feasible)
%   \cmLdInAllQ: a step to the former headline (no headline, no live-feasible)
%   \cmLdInAllT: a step to the former headline (no headline, no live-feasible)
%   \cmLdInAllDM: a step to the former headline (no headline, no live-feasible)
%   \cmLdInBase: a step to the former headline (no headline, no live-feasible)
%   \cmLdInBaseQ: a step to the former headline (no headline, no live-feasible)
%   \cmLdInBaseT: a step to the former headline (no headline, no live-feasible)
%   \cmLdInBaseDM: a step to the former headline (no headline, no live-feasible)
%   \cmLdCMaskLstm: an LSTM step of the campaign (LSTMs parked)
%   \cmLdCMaskLstmQ: an LSTM step of the campaign (LSTMs parked)
%   \cmLdCMaskLstmT: an LSTM step of the campaign (LSTMs parked)
%   \cmLdCMaskLstmDM: an LSTM step of the campaign (LSTMs parked)
%   \cmLdCCadLstm: an LSTM step of the campaign (LSTMs parked)
%   \cmLdCCadLstmQ: an LSTM step of the campaign (LSTMs parked)
%   \cmLdCCadLstmT: an LSTM step of the campaign (LSTMs parked)
%   \cmLdCCadLstmDM: an LSTM step of the campaign (LSTMs parked)
%   \cmLdCSeq: an LSTM step of the campaign (LSTMs parked)
%   \cmLdCSeqQ: an LSTM step of the campaign (LSTMs parked)
%   \cmLdCSeqT: an LSTM step of the campaign (LSTMs parked)
%   \cmLdCSeqDM: an LSTM step of the campaign (LSTMs parked)
%   \cmPNewLive: live_feasible input set (live-feasible removed from the paper)
%   \cmPOldLive: live_feasible input set (live-feasible removed from the paper)
%   \cmKeptLive: live_feasible input set (live-feasible removed from the paper)
%   \cmKeptLiveMin: live_feasible input set (live-feasible removed from the paper)
%   \cmKeptLiveMax: live_feasible input set (live-feasible removed from the paper)
%   \cmLstmCadLive: live_feasible input set (live-feasible removed from the paper)
%   \cmLstmCadBase: LSTM (LSTMs parked)
%   \cmLstmCadAll: LSTM (LSTMs parked)
%   \cmLstmCadBelow: LSTM (LSTMs parked)
%   \cmLstmCadBelowSig: LSTM (LSTMs parked)
%   \cmLstmVsRidgeMin: LSTM (LSTMs parked)
%   \cmLstmVsRidgeMax: LSTM (LSTMs parked)
%   \cmIntraVsRidgeMin: LSTM (LSTMs parked)
%   \cmIntraVsRidgeMax: LSTM (LSTMs parked)
%   \cmIntraAllQ: LSTM (LSTMs parked)
%   \cmIntraAllDM: LSTM (LSTMs parked)
%   \cmNlBestVsH: against the former headline (no headline model)
%   \cmNlBestVsHT: against the former headline (no headline model)
%   \cmNlVsHAbove: against the former headline (no headline model)
%   \cmNlVsHAboveSig: against the former headline (no headline model)
%   \cmPnlTotal: D.14, the former headline's P&L decomposition (parked, user decision)
%   \cmPnlTotalX: D.14, the former headline's P&L decomposition (parked, user decision)
%   \cmPnlTopTenShare: D.14, the former headline's P&L decomposition (parked, user decision)
%   \cmPnlTopTwentyShare: D.14, the former headline's P&L decomposition (parked, user decision)
%   \cmPnlTopTwentyShareX: D.14, the former headline's P&L decomposition (parked, user decision)
%   \cmPnlExTwenty: D.14, the former headline's P&L decomposition (parked, user decision)
%   \cmPnlExTwentySh: D.14, the former headline's P&L decomposition (parked, user decision)
%   \cmPnlMktTwenty: D.14, the former headline's P&L decomposition (parked, user decision)
%   \cmPnlMktTwentyExp: D.14, the former headline's P&L decomposition (parked, user decision)
%   \cmPnlMktTwentyP: D.14, the former headline's P&L decomposition (parked, user decision)
%   \cmPnlDaysAll: D.14, the former headline's P&L decomposition (parked, user decision)
%   \cmPnlDaysAllPct: D.14, the former headline's P&L decomposition (parked, user decision)
%   \cmPnlMEn: D.14, the former headline's P&L decomposition (parked, user decision)
%   \cmPnlMEbuys: D.14, the former headline's P&L decomposition (parked, user decision)
%   \cmPnlME: D.14, the former headline's P&L decomposition (parked, user decision)
%   \cmPnlMEShort: D.14, the former headline's P&L decomposition (parked, user decision)
%   \cmPnlNotMEn: D.14, the former headline's P&L decomposition (parked, user decision)
%   \cmPnlNotME: D.14, the former headline's P&L decomposition (parked, user decision)
%   \cmPnlDiffBuyShare: D.14, the former headline's P&L decomposition (parked, user decision)
%   \cmPnlDiff: D.14, the former headline's P&L decomposition (parked, user decision)
%   \cmPnlDiffX: D.14, the former headline's P&L decomposition (parked, user decision)
%   \cmPnlDiffExTen: D.14, the former headline's P&L decomposition (parked, user decision)
%   \cmPnlDiffHighVix: D.14, the former headline's P&L decomposition (parked, user decision)
%   \cmPnlHighVixN: D.14, the former headline's P&L decomposition (parked, user decision)
%   \cmSvAbove: strategy variations under the former headline (no headline model)
%   \cmSvAboveSig: strategy variations under the former headline (no headline model)
%   \cmSvBestMid: strategy variations under the former headline (no headline model)
%   \cmSvBestMidSh: strategy variations under the former headline (no headline model)
%   \cmSvBestMidT: strategy variations under the former headline (no headline model)
%   \cmSvWingN: strategy variations under the former headline (no headline model)
%   \cmSvWingMid: strategy variations under the former headline (no headline model)
%   \cmSvWingX: strategy variations under the former headline (no headline model)
%   \cmSvWingMidSig: strategy variations under the former headline (no headline model)
%   \cmSvWingXSig: strategy variations under the former headline (no headline model)
%   \cmPtFiredRidgeLean: 2026-10-08 parked (lean set; replaced by an MCS)
%   \cmPtPenRidgeLean: 2026-10-08 parked (lean set; replaced by an MCS)
%   \cmPtFiredLassoLean: 2026-10-08 parked (lean set; replaced by an MCS)
%   \cmPtPenLassoLean: 2026-10-08 parked (lean set; replaced by an MCS)
%   \cmPtFiredEnetLean: 2026-10-08 parked (lean set; replaced by an MCS)
%   \cmPtPenEnetLean: 2026-10-08 parked (lean set; replaced by an MCS)
%
% statements of the text checked on these numbers (2026-10-06; not asserted):
%   holds   5.4: sign(s) lifts every one of the paper's eight forecasts above always short (mid)
%   holds   5.4: R trades above always short (cmRefVsShort > 0)
%   holds   5.4: the FOMC columns' Sharpe gap is not significant (HAC t)
%   holds   4.5: every forecast of the model set leans short (sells on > half the days)
%   holds   5.4: the model set's top Sharpe (subpt_enet_all_features) is not significantly above R (HAC t)
%   holds   5.4: no step of Table 8 moves the Sharpe ratio significantly (HAC t; 0 do)
%   holds   5.4: daily refits make LightGBM more accurate (cmLdCCadLgbmQ < 0)
%   holds   5.4: ... without moving its trade significantly (HAC t)
%   holds   5.4: tuning does not help (TuneRs: QLIKE % > 0)
%   holds   5.4: tuning does not help (TuneOp: QLIKE % > 0)
%   holds   5.4: the de-dup moves the per-bar ridge's QLIKE only at rounding
%   holds   5.4: the de-dup leaves the trade
%   holds   5.4: no position of the per-bar ridge changed
%   FLIPPED 5.4: the specification step moves little on the trade (cmLdSpec)
%   holds   5.4: per-bar beats pooled on QLIKE on HAR + calendar (cmLdPerBarBaseQ < 0)
%   holds   5.4: ... significantly (DM)
%   holds   5.4: ... but not on the trade (HAC t not significant)
%   holds   5.4: no linear forecast's trade moved significantly with the de-dup (HAC t)
%   holds   5.4: the mask made no tree significantly better (DM / HAC t)
%   FLIPPED 5.4: no tree more accurate than the ridge (2 QLIKE differences below 0)
%   FLIPPED 5.4: no tree trades above the ridge (7 Sharpe differences above 0)
%   FLIPPED 5.4: daily refits lower every tree's QLIKE
%   holds   5.4: daily refits move XGBoost on all_features' trade the most
%   FLIPPED 5.4 (was): that move is resolved (HAC t significant)
%   holds   5.4: random search is never more accurate than the shipped settings (0 below)
%   holds   5.4: Optuna is never more accurate than the shipped settings (0 below)
%   holds   5.4: no tuned tree's trade differs significantly from the shipped one (0)
%   FLIPPED 5.4: tuning more often is never less accurate (6 above)
%   holds   5.4: ... and never significantly different on the trade (HAC t)
%   FLIPPED 5.4: the best of 50 trials is never more accurate than of 10 / 25 (7 below)
%   FLIPPED 5.4: at crossed fills the straddle ranks first for every forecast of the model set
%   holds   5.4: wings cost at every width at crossed fills, for every forecast
%   holds   5.4: for every forecast the 20 best days are buys (the tail-trade reading)
%   FLIPPED 4.2: the B contamination moves no Sharpe ratio by 0.1 (model set)
%   FLIPPED 4.5: dropping the censored days moves no Sharpe ratio by more than 0.06
%   holds   4.x: the margin rule lets few retunes fire
%   FLIPPED methods: pre-tuning improves the pooled models' QLIKE (mean change < 0)
\fi

\section{Conclusion}
\label{sec:rh_conclusion}
\begin{itemize}
\item \textbf{Exogenous features lower the QLIKE of the intraday forecast.} Against the HAR benchmark, the
  full predictor set lowers average QLIKE in all \rhrAfPoolPairs{} intraday classes,
  \rhrAfPoolLowerSig{} of the reductions significant at the 5\% level
  (Section~\ref{sec:rh_results}). The tree ensembles have the largest reductions. Among the linear
  models, the elastic net's reduction is significant.
\item \textbf{Squared error does not confirm the QLIKE reductions.} The full set raises the MSE in
  \rhrAfPoolMseHigher{} of the \rhrAfPoolPairs{} classes. No squared-error difference is significant.
\item \textbf{The intraday forecasts remain less accurate than the daily HAR benchmark.} Every intraday
  forecast has a higher average QLIKE than the benchmark forecast, the daily OLS on the HAR benchmark.
  The full predictor set narrows that gap in every class. The linear models' gap is not significant at
  the 5\% level; the trees' gap is (Section~\ref{sec:rh_results_classes}).
\item \textbf{Training on the closing interval is a refinement.} Estimating each model on the last 30
  minutes lowers average QLIKE further, and the full predictor set still lowers QLIKE in every class
  (Section~\ref{sec:rh_econ_last30}).
\item \textbf{The straddle reads the same forecast.} On the option-data subsample, the full predictor set
  raises the sign strategy's Sharpe ratio in every intraday class. Average QLIKE and that Sharpe ratio
  are weakly related (Section~\ref{sec:rh_econ_results}).
\item \textbf{Limits.}
  \begin{itemize}
  \item The accuracy comparison uses $T_{\mathrm{oos}} = \rhoosDays$ trading days, \rhoosSpanStart{} to
    \rhoosSpanEnd.
  \item One index, and one scored interval, the closing half hour.
  \item The straddle is held to expiration with no delta hedge, and trading at the quoted bid and ask
    lowers the Sharpe ratio (Section~\ref{sec:rh_econ}).
  \end{itemize}
\item \textbf{Left open.} Expirations after the evaluation span, a second market, and forecast origins
  earlier in the day.
\end{itemize}
% 2026-10-16 (user): the conclusion now follows the intraday exogenous-feature results.
% The MinGRU bullet and the closing-interval-versus-intraday headline are parked with the old list.
% \begin{itemize}
% \item \textbf{Exogenous predictors add to the HAR benchmark, consistently but modestly.}
%   \begin{itemize}
%   \item $\overline{\QLIKE}$ falls in most pairs and $\SR$ rises in nearly all.
%   \item Few gains are individually statistically significant over $T_{\mathrm{oos}} = \cmNdays$ trading days.
%   \end{itemize}
% \item \textbf{The closing-interval model is more accurate than the intraday model,} significantly so on the
%   HAR benchmark.
%   \begin{itemize}
%   \item The difference in the straddle strategy's $\SR$ is not statistically significant.
%   \end{itemize}
% \item \textbf{A lower QLIKE does not imply a higher Sharpe ratio.} The strategy uses $\hat\sigma^2_t$ only through
%   the side of implied variance on which it falls:
%   \[
%     q_t = \sign\bigl(\hat\sigma^2_t - \mathrm{IV}_t^2\bigr).
%   \]
% \item \textbf{The straddle strategy measures economic value; it is not a claim of implementable profits.}
%   \begin{itemize}
%   \item Every forecast earns a higher $\SR$ than the unconditional short-straddle benchmark ($q_t \equiv -1$).
%   \item Its profits are concentrated in a small number of large-move days.
%   \item At quoted bid/ask prices, rather than mid-quotes, the $\SR$ is materially lower.
%   \end{itemize}
% \item \textbf{A recurrent network (MinGRU) matches the linear models} once regularized like them:
%   \begin{itemize}
%   \item initialized at the ridge solution;
%   \item exogenous predictors shrunk more heavily than the HAR benchmark block, through weight decay in proportion to
%     the linear penalty weights $\omega_j$;
%   % 2026-10-15 (consolidation): parked -- row index s (the periodicity factor's letter) and the estimation sample
%   % W^est_t (the forecast is estimated on the full window W_t; Methods, FLAG (1)). Was:
%   % $\min_{\bm\theta} \sum_{s \in \mathcal W^{\mathrm{est}}_t} \QLIKE(\tilde y_s^{\,2}, \hat y_s^{\,2}(\bm\theta))$, with
%   % $\tilde y_s$ the winsorized target and $s$ the trading days of the estimation sample.
%   \item trained and selected on QLIKE, the evaluation loss,
%     $\min_{\bm\theta} \sum_{t' \in \mathcal W_t} \QLIKE\bigl(\tilde y_{t'}^{\,2},\, \hat y_{t'}^{\,2}(\bm\theta)\bigr)$
%     (up to the floor and clamp of Section~\ref{sec:rh_methods_models}), with $\tilde y_{t'}$ the winsorized target
%     and $t'$ the trading days of the estimation window; QLIKE is scale-free, so this is the same loss on the variance
%     scale.
%   \end{itemize}
%   % 2026-10-15 (notation pass): parked -- the network trains on the winsorized target (fit only, guide section 2),
%   % while the paper now scores against the actual RV; the formula above states the training loss as fitted.
%   % \item trained on the scored loss, $\min \sum_t \QLIKE(RV_t, \hat y_t^2 B_t)$.
% \item \textbf{Limitations:}
%   \begin{itemize}
%   \item $T_{\mathrm{oos}} = \cmNdays$ trading days;
%   \item one index and one intraday interval;
%   \item no delta hedging and no price-impact model;
%   \item tree-ensemble forecasts depend on the CPU instruction set and are not bit-for-bit reproducible across
%     machines.
%   \end{itemize}
% \item \textbf{Future work:}
%   \begin{itemize}
%   \item evaluate the expirations after April 2024 out of sample;
%   \item density forecasts of closing-interval variance;
%   \item a second market;
%   \item forecast origins, and so straddle entries, earlier in the trading day;
%   \item transaction costs calibrated to executed prices.
%   \end{itemize}
% \end{itemize}
